\chapterhead{47}{ORR-8}{Cyber-resilience: Range of Practices}{1}{4}

\lohead{47}{a}{Define cyber resilience and compare recent regulatory initiatives
in the area of cyber resilience.}

\orsub{Cyber-resilience standards and guidelines}

\begin{ordefbox}{Cyber resilience}
Cyber resilience is the ability of an organisation to withstand and recover from
cyberattacks. It is about having a plan in place to minimise the impact of such
attacks and ensure business continuity.
\end{ordefbox}

Jurisdictions worldwide enforce IT and operational risk standards to enhance
cyber resilience, emphasising governance, risk management, IT recovery, and safe
outsourcing practices.

\noindent\begin{minipage}{\textwidth}
\renewcommand{\arraystretch}{1.4}
\begin{tabularx}{\textwidth}{@{}L{0.22\textwidth}Y@{}}
\rowcolor{navy} \thd{Jurisdiction} & \thd{Emphasis}\\
{\tabfont \textbf{1. Hong Kong}} & {\tabfont Targets comprehensive cybersecurity risk management strategies.}\\
\rowcolor{paleblue}
{\tabfont \textbf{2. Singapore}} & {\tabfont Prioritises the detection and mitigation of cyberattacks.}\\
{\tabfont \textbf{3. Brazil}} & {\tabfont Focuses on the formulation and implementation of cybersecurity policies.}\\
\rowcolor{paleblue}
{\tabfont \textbf{4. European Union}} & {\tabfont Develops common procedures through the European Banking Authority for addressing cyber risks.}\\
\end{tabularx}
\end{minipage}
\orfigcap{Regulatory initiatives in cyber resilience, by jurisdiction.}

% ---------------------------------------------------------------------
\lohead{47}{b}{Describe current practices by banks and supervisors in the
governance of a cyber-risk-management framework, including roles and
responsibilities.}

\orsub{Cyber-governance}

\orsubb{1) Cybersecurity strategy}
\begin{lettered}
\item Regulators generally expect institutions to have a cybersecurity strategy
aligned with overall risk management, even if not explicitly mandated.
\item Some jurisdictions require institutions to have a risk management framework
that covers cyber risk.
\item There are three approaches to strategy: regulator-defined,
institution-defined (reviewed by regulators), or regulator review of existing IT
strategies for cyber risk coverage.
\end{lettered}

\orsubb{2) Management roles and responsibilities}
\begin{lettered}
\item The board and senior management are crucial for addressing cyber risk.
\item The three lines of defense model (frontline, risk oversight, independent
assurance) is a common framework, but regulators expect institutions to define
specific responsibilities within this model.
\item Clear and well-defined risk management frameworks are expected.
\end{lettered}

\orsubb{3) Awareness}
\begin{lettered}
\item Regulators encourage a common risk culture and proper employee training.
\item Background checks and nondisclosure agreements might be required for new
hires, with periodic reviews for existing staff.
\end{lettered}

\orsubb{4) Architecture and standards}
\begin{lettered}
\item No consistent regulations for cyber architecture.
\item Focus: up-to-date content storage and periodic self-assessments.
\end{lettered}

\orsubb{5) Workforce}
\begin{itemize}
\item Skills and competencies vary greatly.
\item Standards focus on training the IT workforce to manage risks.
\end{itemize}

% ---------------------------------------------------------------------
\lohead{47}{c}{Explain methods for supervising cyber resilience, testing and
incident response approaches, and cybersecurity and resilience metrics.}

\orsub{Approaches to cyber risk management, incident response and recovery}

\orsubb{1) Supervision of cyber-resilience}

The primary focus of supervision is on the complexity and business model of key
risks. Cybersecurity reviews by regulators may be triggered by institutions' own
risk assessments, by findings from inspections, or by analyses of cyber
incidents. However, there are significant differences in jurisdictional
approaches. Jurisdictions are engaging with industry, not only to address cyber
risk, but to influence behaviour and aid regulatory work.

\orsubb{2) Information security controls}

This focuses on the technical measures institutions take to protect themselves.
It highlights the importance of:

\begin{lettered}
\item Mapping critical business assets and services.
\item Detecting threats through monitoring and surveillance.
\item Implementing security controls for hardware and software.
\end{lettered}

\orsubb{3) Incident response and recovery}

How institutions prepare for and respond to cyber incidents. Key points include:

\begin{lettered}
\item Establishing frameworks for risk detection, response, and recovery.
\item Analysing threats and sharing information to reduce reporting delays.
\item Conducting joint public-private training exercises to improve response
coordination.
\end{lettered}

\orsubb{4) Cybersecurity and resilience metrics}

This addresses the challenge of measuring cyber-resilience. It highlights the use
of:

\begin{lettered}
\item Surveys, incident reports, tests, and inspections for assessment.
\item Management information and indicators provided by institutions.
\end{lettered}

% ---------------------------------------------------------------------
\lohead{47}{d}{Explain and assess current practices for the sharing of
cybersecurity information between different types of institutions.}

\orsub{Sharing cybersecurity information between institutions}

\begin{center}
\begin{tikzpicture}[font=\fontsize{8}{10}\selectfont,
  ent/.style={rounded corners=3pt,inner sep=6pt,text width=2.5cm,align=center,
              text=white,font=\sffamily\bfseries\fontsize{8.4}{10}\selectfont},
  fl/.style={-{Stealth[length=5pt]},line width=1.1pt}]

\node[ent,fill=navy]   (banks) at (-4.6,0)  {Banks};
\node[ent,fill=teal]   (reg)   at ( 1.6,0)  {Regulators};
\node[ent,fill=forest] (sec)   at ( 1.6,-2.5) {Security\\agencies};

\draw[fl,navy] (banks) to[out=130,in=50,looseness=8] (banks)
  node[above=22pt,text=navy,font=\sffamily\fontsize{7.2}{9}\selectfont]
  {1. among banks};
\draw[fl,teal] (banks) -- (reg)
  node[midway,above,text=teal,font=\sffamily\fontsize{7.2}{9}\selectfont]
  {2. banks $\rightarrow$ regulators};
\draw[fl,teal] (reg) to[out=130,in=50,looseness=8] (reg)
  node[above=22pt,text=teal,font=\sffamily\fontsize{7.2}{9}\selectfont]
  {3. among regulators};
\draw[fl,rust] ($(reg.south)+(-0.3,0)$) to[out=235,in=-55]
  ($(banks.south)+(0.3,0)$)
  node[midway,below=3pt,text=rust,font=\sffamily\fontsize{7.2}{9}\selectfont]
  {4. regulators $\rightarrow$ banks};
\draw[fl,forest] (banks) -- (sec)
  node[midway,below,sloped,text=forest,font=\sffamily\fontsize{7.2}{9}\selectfont]
  {5. with security agencies};
\draw[fl,forest] (reg) -- (sec);
\end{tikzpicture}
\end{center}
\orfigcap{The five information-sharing channels and their directions.}

\orsubb{1) Sharing among banks}
\begin{itemize}
\item \textbf{Purpose:} to enhance each bank's preparedness and response to cyber
threats.
\item \textbf{Mechanism:} mostly voluntary, facilitated by regulators or through
specific forums or centres.
\item \textbf{Regulations:} mandatory in some jurisdictions (e.g., Brazil, Japan,
Saudi Arabia).
\item \textbf{Challenges:} data protection laws, trust issues among banks.
\end{itemize}

\orsubb{2) Sharing from banks to regulators}
\begin{itemize}
\item \textbf{Purpose:} to monitor systemic risk, enhance regulatory frameworks,
and improve oversight.
\item \textbf{Requirements:} often mandated by regulations; includes reporting
operational and cyber incidents.
\item \textbf{Variations:} reporting requirements vary in taxonomy, timeframe,
templates, and thresholds.
\end{itemize}

\orsubb{3) Sharing among regulators}

Exchange of information between regulators is less frequent but becoming
increasingly important due to the evolving nature of cyber threats. This can
involve sharing regulatory actions and responses among domestic and cross-border
regulators.

\orsubb{4) Sharing from regulators to banks}

Not widespread; some jurisdictions offer guidance or standards for sharing. Often
anonymised or summarised to protect sensitive information. Selective guidance in
places like the US, UK, China, Singapore and Australia. Mandatory in China;
sharing in most jurisdictions is voluntary.

\orsubb{5) Sharing with security agencies}

Benefits banks, regulators, and the public. Sharing can be voluntary or mandatory
depending on the jurisdiction.

% ---------------------------------------------------------------------
\lohead{47}{e}{Describe practices for the governance of risks of interconnected
third party service providers.}

\begin{orkeybox}
{\sffamily\bfseries\fontsize{8.6}{10}\selectfont\color{navy}Risks of third-party
service providers}\par
\vspace{4pt}
Banks lack control and visibility over third-party systems they rely on.
Regulations require a framework for managing third-party risks. Risks include
security, business continuity, and others.
\end{orkeybox}

\orsubb{1) Governance}
\begin{itemize}
\item Establish a framework defining roles, responsibilities, risk analysis,
monitoring, and risk assessment for outsourcing activities.
\item Risk analysis should cover strategic, compliance, security, business
continuity, vendor lock-in, counterparty, and access risks.
\item Supervisors expect institutions to identify and prioritise third-party
connections and classify incident responses.
\end{itemize}

\orsubb{2) Business continuity and availability}
\begin{itemize}
\item Analyse risks and implement plans to mitigate cyber attack risks from third
parties, and align critical supplier policies with the institution's own
policies.
\item Define recovery plans for critical business activities impacted by
third-party disruptions.
\item Proactively test processes and measures, and conduct scenario analyses.
\end{itemize}

\orsubb{3) Information confidentiality and integrity}
\begin{itemize}
\item Data protection requirements ensure confidentiality and integrity through
confidentiality agreements and security measures. Institutions need to verify,
assess, and monitor third-party security processes.
\item Regulations may specify data location, segregation, use limitations, access
restrictions, and exit strategies.
\end{itemize}

\orsubb{4) Visibility of third-party interconnections}
\begin{itemize}
\item Disclose details about outsourcing arrangements and maintain an inventory
of outsourced functions.
\item Obtain reports from third parties and potentially gain visibility into
sub-outsourcing activities.
\item Some jurisdictions require identification of hardware, software, and
information flows involved.
\end{itemize}

\orsubb{5) Auditing and testing}
\begin{itemize}
\item Guarantee the right to audit and inspect third parties (and potentially
subcontractors).
\item Audit opinions may be based on external or internal audits, or pooled
audits from multiple institutions.
\end{itemize}

\orsubb{6) Resources and skills}
\begin{itemize}
\item Institutions need staff with relevant skills and expertise to assess and
manage emerging risks from new technologies used by third parties.
\item Hire qualified staff to adequately manage and monitor outsourcing
arrangements.
\item Develop plans for handling staffing shortages through consultants or
specialists.
\end{itemize}
